% ---- pg-508 （书页 496） ----

在逻辑斯蒂回归分析中, 检验假设 $H_0$ 常用“负二倍对数似然比”
\phantomsection\label{eq:11-76}
\begin{equation}
-2\ln\left(\frac{L_{\max,\mathrm{Reduced}}}{L_{\max}}\right)
\tag{11-76}
\end{equation}
在这里称为“离差”(deviance). 其分布在假设 $H_0$ 成立时渐近收敛到自由度 1 的卡方分布. 对于大的离差值, 拒绝 $H_0$.

另一种检验单个回归项显著性的方法来自王尔德(Wald, 见文献[17]). 对假设 $H_0:\beta_k=0$, 王尔德检验的检验统计量为 $Z=\hat{\beta}_k/\mathrm{SE}(\hat{\beta}_k)$ 或自由度 1 的卡方分布形式 $Z^2$. 似然比检验优于王尔德检验, 因为似然比检验的显著性水平更接近于名义水平 $\alpha$.

一般的, 如果零假设中有 $m$ 个参数同时为 0, 对于缩减的模型, 离差的渐近分布为自由度 $m$ 的卡方分布.

对于单个的二值观测 $Y_j$, 残差
\[
\frac{Y_j-\hat{p}(\mathbf{z}_j)}{\sqrt{\hat{p}(\mathbf{z}_j)(1-\hat{p}(\mathbf{z}_j))}}
\]
只能取两个值, 因此不是特别有用. 较好的方法是将残差合理地分组, 然后对每组计算总残差. 如果在每组中有 $t$ 个残差, 求这些残差之和再除以 $\sqrt{t}$ 可以使得方差比较具有一致性.

以下我们给出逻辑斯蒂回归的细节和模型检验, 并给出其在分类中的应用.

\subsection{分类}

当观测单位属于总体 1 时, 令响应 $Y=1$; 而当观测单位属于总体 2 时, $Y=0$. (选择响应值取 1 和 0 是任意的, 但有方便性. 在例 11.17 中, 我们用 1 和 2.) 一旦建立起逻辑斯蒂回归函数, 并且有一组来自两个总体的训练样本, 我们可以进一步进行分类. 由于先验概率和代价不易结合到分析中去, 故采用以下分类准则:

\qquad 将 $\mathbf{z}$ 归入总体 1, 若估计的机会比大于 1, 即
\[
\frac{\hat{p}(\mathbf{z}_j)}{1-\hat{p}(\mathbf{z}_j)}=\exp\{\hat{\beta}_0+\hat{\beta}_1 z_{j1}+\cdots+\hat{\beta}_r z_{jr}\}>1
\]
等价的, 我们有简单线性判别准则:

\qquad 将 $\mathbf{z}$ 归入总体 1, 若线性判别函数大于 0, 即
\phantomsection\label{eq:11-77}
\begin{equation}
\ln\left(\frac{\hat{p}(\mathbf{z}_j)}{1-\hat{p}(\mathbf{z}_j)}\right)=\hat{\beta}_0+\hat{\beta}_1 z_{j1}+\cdots+\hat{\beta}_r z_{jr}>0
\tag{11-77}
\end{equation}

\noindent\textbf{例 11.17}\quad (鲑鱼数据的逻辑斯蒂回归)

在例 11.8 中我们引入了鲑鱼数据(\hyperref[tab:11.2]{表 11.2}). 在例 11.8 中, 我们在将鲑鱼分成阿拉斯加鲑鱼和加拿大鲑鱼时只依据年轮而忽略了性别. 如果性别数据也参与分析, 有可能得到更好的分类. 显示板 11.2 包含了用 SAS 软件对鲑鱼数据作逻辑斯蒂回归分析的输出. 这里, 对阿拉斯加鲑鱼 $Y=1$, 而对加拿大鲑鱼 $Y=2$. 预测变量(协变量)为性别(雌性 1, 雄性 2), 淡水生长和海水生长. 由 SAS 输出, 对全局零假设, 似然比检验的结果(缩减的模型只包含 $\beta_0$, 见\hyperref[eq:11-77]{式(11-77)})的显著性水平$<0.0001$. 因此在逻辑斯蒂模型中至少要有一个协变量. 在极大似然估计标题下查看每个个别项, 我们看到王尔德检验建议性别是不显著的($p$ 值$=0.7356$), 而淡水生长和海水生长两个协变量是显著的. 因此可以将性别从模型中去除, 此变量对分类没有意义. 逻辑斯蒂回归可以对不包含性别的模型再次分析, 最终得到的函数可依据判别法则(11-77)将鲑鱼分类为阿拉斯加鲑鱼或加拿大鲑鱼.
